% GENERATED FILE. Edit canonical lessons, then run npm run generate:books.
% canonical-source-hash: fnv1a64:ecd0a47d80ab2c07
% canonical-lessons: TE-C19-vayasu

\chapter{Asking Someone's Age}
\label{ch:te-asking-age}

This chapter is generated from the canonical micro-lessons used by Language Ladder.
Each section stays independently resumable and preserves the authored prerequisite order.

\section[\te{నీ} \te{వయసెంత}?]{\te{నీ} \te{వయసెంత}? — matching Kannada's simple possessive exactly}

\label{lesson:TE-C19-vayasu}



\begin{quote}
\textbf{Warm-up.} {[}PAUSE 2s{]} Telugu matches Kannada closely here — same Sanskrit word, same verbless grammar, just fused together by ordinary Telugu vowel sandhi.
\end{quote}

\begin{cousinweb}
\textbf{\te{వయస్సు}} (\emph{vayasu}) = "\textbf{age}" — the exact same Sanskrit \textbf{\dv{वयस्}} (\emph{vayas}, "age, vigor, strength," PIE cousin of Latin \textbf{vīs}) you just met in Kannada.
\end{cousinweb}

\begin{grammarlens}[title={Grammar Lens: fused by sandhi, still verbless}]
\begin{quote}
\textbf{\te{నీ} \te{వయసెంత}?} — "How old are you?" — literally "\textbf{your age how-much}?"
\end{quote}

(\emph{nī vayas-enta?} — from \emph{nī vayasu} + \emph{enta}, "how much," fused by regular vowel sandhi into \emph{vayasenta}.)

\begin{quote}
\textbf{\te{నా} \te{వయస్సు} \te{ఇరవై}.} — "I am twenty years old." — literally "\textbf{my age}
\end{quote}

{[}is{]} \textbf{twenty}." (\emph{nā vayasu iravai.})

Same shape as Kannada: a plain possessive ("your/my age") sitting next to the number, \textbf{no verb at all}. The only wrinkle is surface-level: everyday spoken Telugu often \textbf{fuses} \emph{vayasu} + \emph{enta} into one word, \emph{vayasenta}, the same kind of vowel sandhi you'd see fusing any two Telugu words in fast speech — the underlying grammar is identical to Kannada's, just written closer together.
\end{grammarlens}

\subsection*{Guided Practice}
{[}PAUSE 1s{]}

\begin{itemize}
\raggedright
  \item {[}YOU SAY: "vayasu" — age{]}
  \item {[}YOU SAY: "nī vayasenta?" — how old are you?{]}
  \item {[}YOU SAY: "nā vayasu iravai" — I am twenty, "my age {[}is{]} twenty"{]}
\end{itemize}

\begin{tcolorbox}[breakable,colback=teal!4,colframe=teal!35!black,title={Wrap-up recall}]
{[}PAUSE 3s{]} Is Telugu's \textbf{\te{వయస్సు}} the same Sanskrit word as Kannada's \emph{vayassu}? (\textbf{Yes} — both from \emph{vayas}, PIE cousin of Latin \emph{vīs}.) What happens to \emph{vayasu} + \emph{enta} in "\emph{nī vayasenta}?"? (They \textbf{fuse} by ordinary vowel sandhi.) Does Telugu use a verb to state age, like Kannada? (\textbf{No} — same verbless possessive shape.)
\end{tcolorbox}
